% =====================================================================
%  Climbing Wall -- Engineering Guide (developer handover, ~12 pages)
%  Compile:  latexmk -pdf engineering_guide.tex
%  Images live in ./images/engineering_guide/
% =====================================================================
\documentclass[11pt,a4paper]{article}
\usepackage{climbingwall}
\usetikzlibrary{positioning,arrows.meta}
\usepackage{longtable}
% ragged-right p-columns: justified cells full of \texttt paths stretch badly
\newcolumntype{L}[1]{>{\raggedright\arraybackslash}p{#1}}
\graphicspath{{images/engineering_guide/}}
\setdocname{Engineering Guide}
% one big diagram: pin it where it is discussed rather than let it float away
\renewcommand{\shotplacement}{p}

\begin{document}

\guidetitle{Engineering Guide}%
  {Handover notes for the next developer --- how the system fits together,
   what it assumes, and where it will bite.}

You have inherited a force-measurement dashboard for the SMS Lab climbing wall.
Four three-axis load-cell boards measure what a climber pulls on each hand and
foot; a Python backend reads them over USB; a React frontend turns the readings
into Newtons, charts them live, records sessions to disk, and computes jump
height.

This guide is not an API reference --- \code{ARCHITECTURE.md} in the repository
root already lists every file and every exported function. What follows is the
part that is not written down in the code: why the pieces are arranged the way
they are, which assumptions are load-bearing, what will look like a bug but
isn't, and what to check before you trust a number the dashboard prints.

% ---------------------------------------------------------------
\section{What to read, and in what order}

The repository documents itself in four Markdown files at the root. They are
maintained and current; do not re-derive from the source what is already
written there.

\begin{center}
\begin{tabular}{@{}L{4.6cm}L{6.3cm}L{4.6cm}@{}}
\toprule
\textbf{File} & \textbf{Covers} & \textbf{Read it} \\
\midrule
\code{README.md} & Setup, running, troubleshooting & Day one \\
This guide & Orientation, invariants, traps, open questions & Day one \\
\code{ARCHITECTURE.md} & Every file, every exported function & When hunting for something \\
\code{measuring\_pipeline.md} & Raw line $\rightarrow$ Newtons $\rightarrow$ storage $\rightarrow$ jump height & Before touching the data flow \\
\code{calibration\_pipeline.md} & The wizard and its maths & Before touching calibration \\
\code{CLAUDE.md} & The same rules, compressed, addressed to AI agents & If you work with one \\
\bottomrule
\end{tabular}
\end{center}

The four LaTeX guides in \code{docs/} (system quick start, user guide,
calibration quick start, calibration guide) are written for the people who
\emph{use} the wall, not for you. Read the user guide once anyway: it is the
fastest description of what the software is supposed to do.

\begin{warnbox}
\textbf{\code{docs/} is gitignored.} The repository's \code{.gitignore} ignores
the whole folder, so every LaTeX source and every built PDF --- including this
guide --- exists only on the machine that made it. If you want the next person
after you to receive it, either un-ignore the sources (keeping the PDFs out, or
in, as you prefer) or move this file into the tracked part of the tree. It is a
one-line change and it is the difference between a handover document and a
local file.
\end{warnbox}

% ---------------------------------------------------------------
\section{The system on one page}

Three things run, and only three:

\begin{enumerate}
  \item \textbf{The boards.} Four PhidgetBridge amplifiers, three channels each
        (X, Y, Z) --- twelve channels total --- plugged into a USB hub, plugged
        into the laptop. There is no Raspberry Pi, no microcontroller firmware,
        and no Web Serial. Older comments that mention any of those are stale;
        trust the code.
  \item \textbf{The backend} (\code{backend/app.py}, Flask, port 5001). It opens
        the twelve channels, holds the latest value of each, and broadcasts them
        over a WebSocket every 10\,ms. It also computes jump height, stores
        recordings, and keeps the calibration files.
  \item \textbf{The frontend} (\code{climbing\_wall\_website/}, React + Vite,
        port 5173). It converts raw readings to Newtons, draws the charts, runs
        the calibration wizard, and owns the zero.
\end{enumerate}

What travels over the WebSocket is deliberately dumb: one line per sample,
twelve comma-separated \emph{raw voltage ratios} of order $10^{-6}$, at roughly
100\,Hz. No units, no calibration, no coordinate transform. Everything that
turns those numbers into physics happens in the browser, in
\code{src/utils/serialParser.ts}.

\begin{infobox}
\textbf{A consequence worth internalising.} Recording CSVs contain
\emph{calibrated Newtons}, not raw ratios --- the conversion happened in the
browser before the data was ever sent to be saved. If the calibration was wrong
when a session was captured, that session is permanently wrong; there is no raw
archive to re-derive it from. Check the calibration before a measurement
campaign, not after.
\end{infobox}

\screenshot{system_architecture}{0.88}{The full data flow, board serial numbers
through to charts. Generated from \code{docs/diagrams/system\_architecture.drawio};
its per-board serial-to-hold mapping matches the live
\code{backend/phidget\_config.json}, \emph{not} the \code{DEFAULT\_CONFIG}
fallback in \code{phidget\_stream.py}, which lists the same serials in a
different order.}

% ---------------------------------------------------------------
\section{Getting it running}

Two terminals, always.

\begin{quote}\ttfamily\small
>> cd backend\\
>> source venv/bin/activate\hfill\textrm{\footnotesize(Windows: \texttt{venv\textbackslash Scripts\textbackslash activate})}\\
>> python app.py\hfill\textrm{\footnotesize$\rightarrow$ \texttt{Running on http://127.0.0.1:5001}}
\end{quote}

\begin{quote}\ttfamily\small
>> cd climbing\_wall\_website\\
>> npm run dev\hfill\textrm{\footnotesize$\rightarrow$ \texttt{http://localhost:5173}}
\end{quote}

Vite proxies \code{/api} and the \code{/api/stream} WebSocket to port 5001, so
the browser only ever talks to 5173. Two details in \code{vite.config.ts} are
load-bearing and have both been debugged the hard way already: the stream rule
must come \emph{before} the generic \code{/api} rule (otherwise the WebSocket
upgrade is caught by the plain HTTP rule), and the target is spelled
\code{127.0.0.1}, not \code{localhost} (Node resolves \code{localhost} to
\code{::1} first, Flask's dev server binds IPv4 only, and the proxy gets
\code{ECONNREFUSED}).

\subsection{Working without hardware}

With no boards attached the app falls back to \textbf{demo mode}: press
\textsf{Start Recording} and \code{useSensorData} generates synthetic values.
This is enough to develop nearly everything --- layout, charts, recordings,
comparison, exports. It cannot do three things, and none of them is a bug:

\begin{itemize}
  \item \textbf{Zero Sensors stays disabled.} Mock data bypasses
        \code{parseSerialLine} entirely, so the tare's sample ring is never fed.
  \item \textbf{The calibration wizard cannot capture.} Same reason; a capture
        times out after 8\,s with no samples.
  \item \textbf{Jump height is fake.} Mock mode returns a random 20--120\,cm
        rather than calling the backend.
\end{itemize}

\subsection{With hardware}

\begin{itemize}
  \item The Phidget driver must be approved once per machine (macOS:
        \textsf{System Settings $\rightarrow$ Privacy \& Security}; Windows:
        install the Phidget22 drivers). Without it, no channel ever attaches.
  \item \textbf{Exactly one program may hold the boards.} They are single-open.
        A second \code{app.py}, or the Phidget Control Panel left running, and
        nothing attaches --- while the backend log looks perfectly healthy.
  \item Restarting the backend requires re-clicking \textsf{Connect Device} in
        the browser. The socket is gone; the UI does not reconnect itself.
  \item Boards may be plugged in before or after the backend starts. Channels
        are opened non-blocking and attach whenever hardware appears, usually
        within a second or two.
\end{itemize}

% ---------------------------------------------------------------
\section{The one invariant}

\begin{warnbox}
\textbf{Everything is stored in the sensor (board) frame. Never rotate on
ingestion.}
\end{warnbox}

The wall is mounted with a backward decline of $\theta$ degrees from vertical
(16$^\circ$ by default). The load cells are bolted to the wall, so their axes
tilt with it:

\begin{center}
\begin{tabular}{@{}ll@{}}
\toprule
\multicolumn{2}{@{}l}{\textbf{Sensor frame} --- canonical, what everything stores} \\
\midrule
$X$ & in-plane, along the wall, positive up-slope \\
$Y$ & sideways, positive right \\
$Z$ & normal to the wall surface, positive out toward the climber \\
\midrule
\multicolumn{2}{@{}l}{\textbf{World frame} --- derived, display only} \\
\midrule
$X_{\text{world}} = X\cos\theta - Z\sin\theta$ & true vertical, up \\
$Z_{\text{world}} = X\sin\theta + Z\cos\theta$ & true horizontal, out of the wall \\
$Y_{\text{world}} = Y$ & unchanged \\
\bottomrule
\end{tabular}
\end{center}

The live buffer, IndexedDB, and the backend CSV all hold sensor-frame Newtons.
The rotation is applied at \emph{render time} only, by
\code{toDisplayFrameReadings} in \code{wallGeometry.ts}, and only when the user
has explicitly opted into the world view.

$\theta$ legitimately appears in exactly three places:

\begin{enumerate}
  \item \textbf{In the calibration wizard}, to split one hang into its $X$ and
        $Z$ components ($W\cos\theta$ and $W\sin\theta$).
  \item \textbf{In the world-frame display}, as the opt-in rotation above.
  \item \textbf{In the backend jump algorithm}, which projects wall-frame forces
        to global vertical --- see the open question in section 12.
\end{enumerate}

Nowhere else. This invariant was inverted once already and produced
double-corrected jump heights that looked plausible and were wrong. If you find
yourself applying $\theta$ in \code{serialParser.ts}, in \code{useSensorData.ts},
or anywhere before a \code{POST}, stop: that is the bug you are about to
introduce, not the fix you are looking for.

% ---------------------------------------------------------------
\section{Data contracts}

\subsection{Twelve channels, one order}

Every length-12 array in the frontend is in \textbf{GUI-slot order}:

\begin{center}
\code{[Left Hand, Right Hand, Left Foot, Right Foot]} $\times$ \code{(X, Y, Z)},
\quad flat index $=$ \code{board * 3 + axis}
\end{center}

The hardware streams a \emph{different} order (group 0 = Left Foot, 1 = Left
Hand, 2 = Right Foot, 3 = Right Hand), and \code{SENSOR\_GROUP\_ORDER = [1, 3,
0, 2]} in \code{serialParser.ts} remaps it. Which physical board is which group
is not in the code at all --- it is \code{bridgeSerials} in
\code{backend/phidget\_config.json}, where the list index is the group number.
\textbf{If sensors appear swapped in the GUI, reorder that list.} Nobody should
ever need to edit \code{SENSOR\_GROUP\_ORDER}.

Recording CSVs follow the GUI order too: \code{S1} is Left Hand, \code{S2} Right
Hand, \code{S3} Left Foot, \code{S4} Right Foot.

\subsection{Two representations of a sample}

\begin{center}
\begin{tabular}{@{}L{2.4cm}L{4.1cm}L{2.2cm}L{6.4cm}@{}}
\toprule
\textbf{Name} & \textbf{What it is} & \textbf{Magnitude} & \textbf{Who consumes it} \\
\midrule
\code{signedRaw} & Post axis-sign, \emph{pre}-tare, \emph{pre}-scale & $\sim10^{-6}$ & The calibration wizard and the zero \\
\code{values} & Fully calibrated Newtons, sensor frame & $\sim10^{2}$ & Everything else: buffer, charts, IndexedDB, CSV, jump test \\
\bottomrule
\end{tabular}
\end{center}

The split is the reason the wizard and the tare work at all: both average
\code{signedRaw}, so their measurements do not depend on whatever offset or
scale happens to be active while they run. Feeding them \code{values} would make
each calibration depend on the previous one. The transform itself is a fixed
four-step chain:

\begin{figure}[H]\centering
\begin{tikzpicture}[
  font=\footnotesize,
  box/.style={draw=line, rounded corners=2pt, fill=soft, inner sep=5pt,
              minimum height=8mm, minimum width=15mm, align=center},
  lbl/.style={align=center, text=ink},
  arr/.style={-{Latex[length=4pt]}, draw=accent, thick}]
  \node[lbl] (in) {12 raw\\ratios};
  \node[box, right=5mm of in]  (r) {remap};
  \node[box, right=5mm of r]   (s) {sign};
  \node[box, right=5mm of s]   (t) {tare};
  \node[box, right=5mm of t]   (c) {scale};
  \node[lbl, right=5mm of c]   (out) {Newtons\\(sensor frame)};
  \draw[arr] (in) -- (r); \draw[arr] (r) -- (s); \draw[arr] (s) -- (t);
  \draw[arr] (t) -- (c);  \draw[arr] (c) -- (out);
  \node[lbl, below=8mm of s, text=muted]
    (tap) {\code{signedRaw} $\rightarrow$ calibration wizard, Zero Sensors};
  \draw[arr, draw=muted] (s.south) -- (tap.north);
\end{tikzpicture}
\caption{\code{parseSerialLine} in \code{serialParser.ts}. The tap after the
sign step is what keeps calibration and tare independent of each other.}
\end{figure}

\subsection{Sample numbers are not indices}

Each buffered reading carries an absolute, monotonically increasing
\code{sampleNumber}. The buffer is trimmed when it exceeds 360\,000 samples, so
array indices shift underneath you; sample numbers do not. The jump test marks
its window by sample number for exactly this reason, and \textsf{Start Fresh}
resets the numbering, which makes a stale jump window unreachable rather than
silently wrong.

The backend's recording window (\code{?from=} and \code{?to=} on the data
endpoint) is the opposite: those are 0-based \emph{row positions} in the file,
not sample numbers. The docstring on \code{\_slice\_and\_downsample} says so;
believe it.

% ---------------------------------------------------------------
\section{Where to change what}

\begin{longtable}{@{}L{5.6cm}L{9.6cm}@{}}
\toprule
\textbf{You want to} & \textbf{Touch} \\
\midrule
\endfirsthead
\toprule
\textbf{You want to} & \textbf{Touch} \\
\midrule
\endhead
Fix a sensor showing up as the wrong body part &
  \code{backend/phidget\_config.json} $\rightarrow$ \code{bridgeSerials}. No code change. \\
Change the sample rate &
  \code{phidget\_config.json} $\rightarrow$ \code{dataIntervalMs}. This is also
  the single source of truth for the jump test's $dt$. \\
Change the amplifier gain &
  \code{phidget\_config.json} $\rightarrow$ \code{bridgeGain}. Invalidates every
  saved scale --- recalibrate afterwards. \\
Change how raw becomes Newtons &
  \code{src/utils/serialParser.ts}. \\
Change the calibration maths &
  \code{src/utils/calibration.ts}; the wizard's UI is
  \code{dashboard/CalibrationModal.tsx}, its capture machinery is
  \code{hooks/useCalibration.ts}. \\
Change the zero &
  \code{src/utils/runtimeOffset.ts} plus \code{hooks/useTareOffset.ts}. \\
Change anything about the wall angle or the world view &
  \code{src/utils/wallGeometry.ts}. Extend \code{wallDeclineSelfCheck()} in the
  same commit. \\
Change buffering, the UI sync rate, or autosave &
  \code{hooks/useSensorData.ts}. \\
Change a chart &
  \code{utils/dataProcessing.ts} (data building), \code{utils/chartOptions.ts}
  (behaviour), \code{components/dashboard/tabs/*} (layout). \\
Add a backend endpoint &
  \code{backend/app.py}, plus a thin client in \code{src/utils/*Api.ts}.
  Components never \code{fetch} directly. \\
Change the jump physics &
  \code{backend/app.py} $\rightarrow$ \code{calculate\_jump\_height\_with\_angle}. \\
Wire a new hook into the app &
  \code{components/sensor-dashboard.tsx} --- and nothing else. \\
\bottomrule
\end{longtable}

\code{sensor-dashboard.tsx} is the only orchestrator in the frontend. It
instantiates the hooks, routes each incoming line to its three sinks through
stable refs (which is how the circular dependency between the stream and its
consumers is broken), and composes the layout. No business logic belongs in it.
When you are tempted to put "just one calculation" there, that calculation
belongs in a hook or a util.

% ---------------------------------------------------------------
\section{Calibration, profiles, and the zero}

Three layers, deliberately separated, and the separation is the design:

\begin{center}
\begin{tabular}{@{}L{3.2cm}L{5.0cm}L{6.8cm}@{}}
\toprule
\textbf{Layer} & \textbf{What} & \textbf{Where it lives} \\
\midrule
Active calibration &
  \code{axisSigns} and \code{axisScales}, 12 numbers each &
  \code{src/config/calibration\_settings.json} (written by the backend),
  mirrored in \code{localStorage["calibration"]} \\
Profile library &
  Named calibrations you can load &
  \code{backend/calibration\_profiles.json} --- gitignored, so a fresh clone has none \\
Zero (tare) &
  A volatile offset in signed-raw space &
  \code{localStorage["tareOffset"]} only. Never saved with the calibration,
  never sent to the backend \\
\bottomrule
\end{tabular}
\end{center}

The zero was part of the saved calibration once and was deliberately pulled out.
What stays in the calibration is only what genuinely belongs to the hardware:
the sign of each axis and its Newtons-per-raw-unit slope. Where zero happens to
sit today is a property of the moment --- temperature, a hold that has been
leaned on, a cable being dressed --- so it is measured live from
\textsf{Zero Sensors} (about 150 samples, roughly 1.5\,s) and cached only so a
page reload does not lose it.

Note the ordering consequence: the tare is applied \emph{before} the scale, in
signed-raw space, which is why it stays valid when you change the scale and why
the wizard is unaffected by whatever zero is in effect.

\begin{infobox}
\textbf{$\theta$ is not part of the calibration.} It is stored separately
(\code{localStorage["cw:wallDeclineDeg"]}) and is only needed \emph{during} a
hang capture, to interpret how much of the hung weight landed on $X$ versus $Z$.
The fitted scales are an intrinsic sensor property; in the ideal orthogonal
model the $\cos\theta$ and $\sin\theta$ factors cancel in the least-squares
slope. Changing the wall angle does not require recalibrating.
\end{infobox}

Editing a profile's numbers by hand is intentionally not offered. The only way
to change a profile's values is to load it, recalibrate against real weights,
and save the changes back --- the \textsf{Recalibrate} button in the profiles
modal does exactly that. If a future request asks for raw number entry, the
answer that was already given once is no.

% ---------------------------------------------------------------
\section{Decisions that look like bugs}

Every item here was a deliberate choice. Each has a failure mode you will
reintroduce if you "clean it up".

\begin{center}
\begin{tabular}{@{}L{5.0cm}L{10.2cm}@{}}
\toprule
\textbf{The odd thing} & \textbf{Why, and what breaks if you undo it} \\
\midrule
The live buffer is a mutable ref, pushed to in place &
  At 100\,Hz, React state would mean thousands of re-renders per second. A
  50\,ms interval copies a display slice instead ($\sim$20\,fps). Never hand
  \code{allSensorDataRef.current} to React state or Chart.js directly. \\
\code{decimate} sometimes returns its input array &
  Same reference means memoised chart components can skip work. But with the
  window set to "all", that input \emph{is} the live buffer, so
  \code{buildDisplaySlice} copies in that case --- otherwise Chart.js would read
  an array being mutated underneath it. \\
Charts are capped at 1000 points, twice &
  Once in the frontend (\code{decimate}), once in the backend
  (\code{MAX\_CHART\_POINTS}). Even-stride selection plus the 5-sample moving
  average is visually faithful and keeps payloads and render cost bounded. \\
The buffer is trimmed in 36\,000-sample chunks &
  Capped at 360\,000 ($\sim$1 hour). Without it, an all-day session grows
  unbounded and the 5\,s IndexedDB autosave re-serialises the lot each time. \\
A slow WebSocket client drops its own samples &
  Each client has a bounded queue; a full queue drops rather than blocks, so one
  stalled browser tab cannot stall the broadcast thread for everyone. \\
The broadcast loop skips when nothing is attached &
  No clients or no attached channels means no line is sent at all --- which is
  why a completely unplugged rig produces silence rather than a stream of zeros. \\
The recordings download endpoint ignores the downsampler &
  The chart endpoint thins deliberately; an export must not. Someone exporting
  "their data" gets the file that was written at save time, every sample. \\
\bottomrule
\end{tabular}
\end{center}

% ---------------------------------------------------------------
\section{Hardware realities}

The physical rig does not behave like an ideal data source, and several
decisions in \code{phidget\_stream.py} exist only because of that.

\begin{itemize}
  \item \textbf{Channels are opened non-blocking.} The backend starts, and runs,
        with no hardware present. Channels attach whenever boards appear. This
        is why the frontend waits 2\,s after the socket opens before judging how
        many channels are present.
  \item \textbf{A detached channel is forced to \code{0.0}}, not left at its last
        reading. A frozen last value is indistinguishable from a real constant
        force in a chart and in a recording; a hard zero at least reads as
        obviously wrong.
  \item \textbf{Partial attachment is normal and is warned about.} The banner
        names the sensor ("Only 9 of 12 sensor channels are attached --- check:
        Right Hand"), and missing channels read 0.
  \item \textbf{Flask runs with \code{use\_reloader=False}.} The reloader forks a
        second process that would also open the channels; single-open means both
        processes then fail to attach and the symptom looks like broken hardware.
  \item \textbf{Network-server discovery is deliberately off.} With the driver
        working, a channel silently attaching through a Phidget Network Server
        instead of USB would be nondeterministic and miserable to debug.
\end{itemize}

\begin{warnbox}
Hardware is never the ideal on paper. Load cells drift, amplifiers have offsets,
a board that has been re-bolted reads differently. The calibration wizard and
the live zero are not ceremony to be optimised away --- they are the knobs that
absorb the difference between the model and the wall.
\end{warnbox}

% ---------------------------------------------------------------
\section{Checking your work}

There is no test framework, and that is a considered choice rather than
neglect. What exists instead:

\begin{enumerate}
  \item \textbf{The build gate.} \code{npm run build} runs \code{tsc} in strict
        mode (\code{noUnusedLocals}, \code{noUnusedParameters}) and then
        bundles. It catches most breakage. This is \emph{the} gate --- run it
        before every commit.
  \item \textbf{\code{wallDeclineSelfCheck()}} in \code{wallGeometry.ts}. Three
        assertions covering the rotation model: $\theta = 0$ is the identity, a
        pure vertical load resolves to all-$X$, a pure normal push splits as
        $(-\sin\theta, \cos\theta)$. It runs automatically in the browser console
        in DEV. \textbf{Extend it when you touch the geometry} rather than
        adding a test runner.
  \item \textbf{\code{backend/test\_recordings.py}}. Plain \code{assert}s, run
        with \code{python test\_recordings.py}. It covers the recording
        window/downsample arithmetic, whose edge cases (reversed bounds,
        out-of-range bounds) are reachable by dragging a range slider.
  \item \textbf{\code{backend/board\_diagnostic.py}}. Not part of the app: a
        forensic tool that answers "is this board's problem calibration, or is it
        mechanical?" from a recording of known-magnitude pulls. Its argument is
        that a calibration fault can only rescale each axis independently, so if
        no per-axis scaling reproduces the known applied magnitudes,
        recalibrating will never fix that board.
  \item \textbf{The manual smoke path}, when hardware is attached: connect,
        confirm 12 of 12 channels, zero with the wall unloaded, hang a known
        weight and check the Newtons are sane, record, save, reload the
        recording from the Recordings tab.
\end{enumerate}

Both automated gates were green on 2026-08-20, with the in-flight
recordings-browser changes described in section 12 present in the tree.

\begin{infobox}
If you add non-trivial logic --- a parser, a loop, a piece of physics --- leave
one runnable check behind it, in the style already used here: an
\code{assert}-based \code{demo()} next to the code, or three lines added to
\code{wallDeclineSelfCheck()}. Do not introduce a test framework to hold six
assertions.
\end{infobox}

% ---------------------------------------------------------------
\section{Traps}

\begin{center}
\begin{tabular}{@{}L{5.4cm}L{9.8cm}@{}}
\toprule
\textbf{Thing} & \textbf{Why it bites} \\
\midrule
\code{allSensorDataRef.current} is mutated in place &
  Handing it to React state means no re-render (same reference) and an array
  changing underneath Chart.js. \\
Mock mode bypasses \code{parseSerialLine} &
  So \textsf{Zero Sensors} stays disabled with no hardware --- the tare ring is
  never fed. Not a bug. \\
Vite ignores \code{calibration\_settings.json} &
  The backend writes \code{src/config/calibration\_settings.json} on every
  calibration save. Without the ignore in \code{vite.config.ts},
  Vite full-page-reloads and drops the stream mid-wizard. \\
Flask runs \code{use\_reloader=False} &
  Two processes cannot both open single-open boards. \\
\code{calibration\_profiles.json} is gitignored &
  The profile library is local-only. A fresh clone starts with none, and nothing
  warns you. \\
The Vite proxy targets \code{127.0.0.1} &
  Not \code{localhost}. IPv6 resolution order versus Flask's IPv4-only bind. \\
IndexedDB is at store version 2 &
  The v1 $\rightarrow$ v2 upgrade \emph{drops} the buffer, because the canonical
  frame changed from world to sensor. Stale data must never be reinterpreted. \\
Recording CSVs hold calibrated Newtons &
  Not raw ratios. A bad calibration at capture time cannot be undone later. \\
\bottomrule
\end{tabular}
\end{center}

% ---------------------------------------------------------------
\section{Open questions and rough edges}

This is the honest part of the handover: things that are unresolved, unverified,
or simply unfinished as of 2026-08-20.

\subsection{The jump projection axes --- verify before trusting absolute heights}

\code{calculate\_jump\_height\_with\_angle} projects the summed hand and foot
forces onto a global vertical axis like this:

\begin{quote}\ttfamily\small
foot\_global\_z = foot[:,2]*cos($\theta$) + foot[:,1]*sin($\theta$)
\end{quote}

Column 2 is $Z$ (wall normal) and column 1 is $Y$ (sideways). But the frontend
sends sensor-frame triples $[X, Y, Z]$, and under this project's own convention
the weight-bearing axes are $X$ (in-plane, $\approx W\cos\theta$) and $Z$
(normal, $\approx W\sin\theta$), while $Y$ is lateral and physically independent
of the wall tilt. A tilt projection in the $X$--$Z$ plane that consumes $Y$ and
ignores $X$ does not obviously correspond to the geometry documented everywhere
else.

The recordings support the convention rather than the code: in
\code{peter\_test\_1\_20260812\_125458.csv} the large mean loads sit on $X$
($-83$, $-135$, $-54$\,N across the three live boards) with $Z$ at $+37$ to
$+49$\,N and $Y$ small, exactly as a hang on a $16^\circ$ overhang should look.

This may still be intentional --- the function's docstring is ambiguous about
whether the height it returns is vertical or wall-perpendicular, and it applies
a further $\cos\theta$ at the end. Do not "fix" it from the armchair.
\textbf{Measure a jump of known height and compare} before changing anything;
then, whichever way it comes out, write the answer down in
\code{measuring\_pipeline.md}. Until then, treat absolute jump heights as
unvalidated --- which is also what the UI says, since the tab is labelled
\textsf{Jump Test (BETA)}.

Two smaller things in the same function: the "first 1 second" comment sits above
\code{total\_force\_z[: still\_end * 3]}, which is three seconds; and
\code{i\_max\_accel} is the argmax of foot \emph{force}, not acceleration.
Neither is necessarily wrong, but both will mislead the next reader.

\subsection{The Right Foot board}

The fourth GUI slot (\code{S4}, Right Foot, bottom right) reads flat zero
through recent captures, and the commit history is blunt about it ("bottom right
board is malicious"). \code{board\_diagnostic.py} exists because this was being
chased. Before blaming code for a missing channel, check the attach count in the
warning banner and run the diagnostic against a known-weight recording.

\subsection{Work in flight}

The working tree on branch \code{minimalstic\_code} carries an unfinished
recordings browser: listing all recordings rather than the newest five
(\code{?limit=N} now opt-in), rename, delete, verbatim CSV download, and a
\code{[from, to)} row window on the data endpoint so zooming actually raises
resolution instead of magnifying the same thinned points. The window arithmetic
is covered by \code{test\_recordings.py}; the UI
(\code{RecordingsBrowserModal.tsx}) is newer than the rest. Finish or discard it
deliberately --- do not let it rot in the tree.

\subsection{Smaller items}

\begin{itemize}
  \item The frontend bundles to a single $\sim$820\,kB chunk. Nobody has needed
        to care yet; it is a lab tool on a local network, not a public site.
        \code{katex} is in there to render the wizard's formulas.
  \item \code{useCalibration.ts}'s header comment still names
        \code{captureAxis} and \code{commitAxis}; the real exports are
        \code{captureAverages} and \code{commitChannels}. Typical drift --- fix
        it as you pass.
  \item \code{docs/} is gitignored, as noted in section 1.
\end{itemize}

% ---------------------------------------------------------------
\section{House rules}

Conventions this codebase actually follows. Match them and your changes will
read like they belong.

\begin{itemize}
  \item \textbf{Comments explain \emph{why}}, never what. A physical ceiling, a
        non-obvious threshold, a trap that has already been hit once. Migration
        history ("this used to be on the Pi") belongs in git, not in a file.
  \item \textbf{No line numbers in documentation.} They rot within a single
        commit. Reference the file and the function name.
  \item \textbf{Prefer configuration over code.} The board mapping, sample rate,
        and gain all live in \code{phidget\_config.json} specifically so that
        the most common field fix requires no developer.
  \item \textbf{Keep the four Markdown docs in step with the code.} They are the
        contract the next person reads; a stale doc is worse than no doc.
  \item \textbf{Smallest change that works.} This codebase is roughly 9\,500
        lines including the backend, and it stays legible because features were
        repeatedly cut rather than added. Before building an abstraction, check
        whether the standard library, the platform, or an already-installed
        dependency covers it.
\end{itemize}

% ---------------------------------------------------------------
\section{Your first hour}

\begin{enumerate}
  \item Start both servers and open the dashboard in demo mode. Press
        \textsf{Start Recording} and watch the charts fill.
  \item Read \code{measuring\_pipeline.md} end to end with the app open beside
        it. It is the one document that will save you the most time.
  \item Run both self-checks: \code{npm run build}, and
        \code{python backend/test\_recordings.py}. Open the browser console and
        confirm \code{wallDeclineSelfCheck()} reports three passes.
  \item Open the Recordings tab and load a saved capture from
        \code{backend/saved\_recordings/} --- real data, real shapes, and a good
        feel for what the numbers should look like.
  \item Read \code{serialParser.ts} in full. It is 84 lines and it is the heart
        of the system.
\end{enumerate}

\vspace{4pt}
{\color{muted}\small Written 2026-08-20, against branch \code{minimalstic\_code}.
When it goes stale, edit it --- it is a source file like any other.}

\end{document}
